%\subsection{Proof Techniques}\label{sec:proof-techniques}
%In this section, we outline our proof techniques for Theorems \ref{thm:complete-splitting-off}, \ref{thm:k-EC-hypergraph-characterization}, \ref{thm:rooted-steiner-orientation-in-graphs}, and \ref{thm:rooted-steiner-orientation-in-hypergraphs}.

\subsection{Proof Technique for Theorem \ref{thm:k-EC-hypergraph-characterization}}
We outline the proof technique for Theorem \ref{thm:k-EC-hypergraph-characterization}. The reverse direction follows by observing that if a hypergraph is $k$-hyperedge-connected, then both operations in the statement of the theorem preserve $k$-hyperedge-connectivity. We sketch a proof of the forward direction. 
The proof is by induction on the number of hyperedges plus the number of vertices. 
First, suppose that there exists a hyperedge $e \in E$ such that $G - e$ is still $k$‐hyperedge‐connected. We note that deleting the hyperedge $e$ is the inverse of operation (1). Consequently, the proof follows by deleting the hyperedge $e$, using the induction hypothesis on the resulting hypergraph $H$, and then noting that the hypergraph $G$ is obtained from $H$ by operation (1). 
Next, suppose that there does not exist a hyperedge $e \in E$ such that $G-e$ is $k$-hyperedge-connected. We call such a hypergraph to be minimally $k$-hyperedge-connected. In \Cref{lem:minimally-k-EC-hypergraph-has-degree-k-vertex}, we show that a minimally $k$-hyperedge-connected hypergraph contains a vertex $u$ with degree exactly $k$. By \Cref{thm:complete-splitting-off}, there exists a global-connectivity preserving complete h-splitting-off at the vertex $u$ from the hypergraph $G$. Let $H$ be a global-connectivity preserving complete h-splitting-off at the vertex $u$ from the hypergraph $G$. 
We note that complete h-splitting-off at $u$ followed by deletion of the vertex $u$ is the inverse of operation (2) at $u$. 
Consequently, the proof follows by using the induction hypothesis on the hypergraph $H-u$ and then noting that the hypergraph $G$ is obtained from $H-u$ by operation (2).

%\shubhang{sketch for forward direction: First, suppose that there exists a hyperedge $e \in E$ such that $H - e$ is still $k$‐hyperedge‐connected. We note that deleting the hyperedge $e$ is the inverse of operation (1). Consequently, the proof follows by induction. Alternatively, suppose that there does not exist a hyperedge $e \in E$ such that $H -e$ is $k$-hyperedge-connected, i.e. $H$ is minimally-k-hyperedge-connected. In \Cref{lem:minimally-k-EC-hypergraph-has-degree-k-vertex}, we show that there exists a vertex $u \in V$ with degree exactly $k$. By \Cref{thm:complete-splitting-off} there exists a global-connectivity preserving complete h-splitting-off of $u$ from the hypergraph $H$. We note that complete h-splitting-off at $u$ is the inverse of operation (2) at $u$. Consequently, the proof follows by induction.}


\subsection{Proof Technique for Theorems \ref{thm:rooted-steiner-orientation-in-graphs} and \ref{thm:rooted-steiner-orientation-in-hypergraphs}}
We outline the proof technique for Theorem \ref{thm:rooted-steiner-orientation-in-hypergraphs} and will remark after the proof about how it also implies a proof for Theorem \ref{thm:rooted-steiner-orientation-in-graphs}. 
Our proof of Theorem \ref{thm:rooted-steiner-orientation-in-hypergraphs} will be in three steps. 
Let us denote the set of non-terminals as \emph{Steiner vertices}. 
Our first step is to obtain a hypergraph $H = (T, E_H)$ by applying our local connectivity preserving complete h-splitting-off at each Steiner vertex of $G$ (sequentially, in arbitrary order of the Steiner vertices) and deleting the isolated  vertices. 
We note that deleting the isolated vertices ensures that the vertex set of $H$ is the set of terminals $T$. 
Moreover, our local connectivity preserving complete h-splitting-off ensures that the hypergraph $H$ is $2k$-hyperedge-connected since the hypergraph $G$ is Steiner $2k$-hyperedge-connected. 
Our second step is to show that this hypergraph $H=(T, E_H)$ admits a rooted $k$-hyperarc-connected orientation. 
A known characterization for the existence of a rooted $k$-hyperarc-connected orientation of a hypergraph is that the hypergraph is \emph{$k$-weak-partition-connected} (see Definition \ref{defn:weak-partition-connectivity} for the definition of weak-partition-connectivity and Theorem \ref{thm:frank:weak-pc-iff-rooted-hyparc-conn-orientation} for the characterization). We mention that the notion of weak-partition-connectivity in hypergraphs has been used recently in the  context of coding theory \cite{GLSTW22, AGL23}. 
In order to use the characterization for the existence of a rooted $k$-hyperarc-connected orientation of a hypergraph, we relate the connectivity of a hypergraph to its \emph{weak-partition-connectivity} and conclude that if $H$ is $2k$-hyperedge-connected, then it is $k$-weak-partition-connected (see Lemma \ref{lem:2k-conn-implies-k-weak-pc}). 
Consequently, the hypergraph $H$ admits a rooted $k$-hyperarc-connected orientation. 
We note that such an orientation of $H$ is equivalent to a Steiner rooted $k$-hyperarc-connected orientation of $H$ since the vertex set of $H$ is the set $T$ of terminals. 
Our third step is to use this Steiner rooted $k$-hyperarc-connected orientation of $H$ to obtain a Steiner rooted $k$-hyperarc-connected orientation of the hypergraph $G$: we will see that there is a natural way to extend the orientation of hyperedges while reversing the h-splitting-off operations to preserve Steiner rooted $k$-hyperarc-connected property (see Lemma \ref{lem:lifting-Steiner-orientation-through-splitting-off}). This would complete the proof of \Cref{thm:rooted-steiner-orientation-in-hypergraphs}. 

We note that if the hypergraph $G$ is a graph, then the same proof above obtains the required graph orientation, thus proving \Cref{thm:rooted-steiner-orientation-in-graphs}.
In particular, to prove Theorem \ref{thm:rooted-steiner-orientation-in-graphs}, we start from a graph $G=(V, E)$ that is Steiner $2k$-edge-connected, but our local connectivity preserving complete h-splitting-off operations at Steiner vertices results in a hypergraph $H=(T, E_H)$ that is $2k$-hyperedge-connected; by Lemma \ref{lem:2k-conn-implies-k-weak-pc}, the hypergraph $H$ is $k$-weak-partition-connected; now Theorem \ref{thm:frank:weak-pc-iff-rooted-hyparc-conn-orientation} gives a rooted $k$-hyperarc-connected orientation of the resulting hypergraph. Such an orientation is extended to a  Steiner rooted $k$-hyperarc-connected orientation of the \emph{graph} $G=(V, E)$ using Lemma \ref{lem:lifting-Steiner-orientation-through-splitting-off}. Essentially, the proof starts from the given graph, obtains a related hypergraph, orients that hypergraph, and extends that orientation of the hypergraph back into a desired orientation of the given graph. 

Our proof technique for Theorems \ref{thm:rooted-steiner-orientation-in-graphs} and \ref{thm:rooted-steiner-orientation-in-hypergraphs} also leads to an alternate proof of Menger's theorem in \emph{undirected} graphs and hypergraphs (edge-disjoint version)---see Section \ref{sec:rooted-steiner-connected-orientations}. 

\subsection{Proof Technique for Theorem \ref{thm:complete-splitting-off}}
We prove a more general statement that implies Theorem \ref{thm:complete-splitting-off}. We begin with the definitions needed for the more general statement. 
%Let $(H=(V, E), w:E\rightarrow \Z_+)$ be a hypergraph. For $X\subseteq V$, let $B_H(X):=\{e\in X: e\cap X\neq \emptyset\}$. We define the coverage function $b_{(H, w)}:2^V\rightarrow \Z_{\ge 0}$ by $b_{(H,w)}(X):=\sum_{e\in B_H(X)}w(e)$ for every $X\subseteq V$. For a vertex $v\in V$, we use $b_{(H, w)}(v)$ to denote $b_{(H, w)}(\{v\})$. We note that $b_{(H, w)}(v)$ denotes the degree of the vertex $v$. 

\begin{definition}
    Let $V$ be a finite set, $p: 2^V\rightarrow \Z$ be a set function,  and $(H=(V, E), w:E\rightarrow \Z_+)$ be a hypergraph. 
    %We will denote the maximum function value of $p$ by $K_p$, i.e., $K_p\coloneqq \max\{p(X): X\subseteq V\}$. 
    %\knote{Is it really worth defining $K_p$ here?}\kristof{We do not use it in this section, I think we could move its definition to the Preliminaries.}\shubhang{Moved to preliminaries.}
\begin{enumerate}
    \item The set function $p$
    \begin{enumerate}
    \item is \emph{symmetric} if $p(X)=p(V-X)$ for every $X\subseteq V$, and 
    \item is \emph{skew-supermodular} if for every $X, Y\subseteq V$, at least one of the following inequalities hold: 
    \begin{enumerate}
    \item $p(X) + p(Y) \leq p(X\cap Y) + p(X\cup Y)$. %If this inequality holds, then we say that $p$ is \emph{locally supermodular} at $X, Y$.
    \item $p(X) + p(Y) \leq p(X - Y) + p(Y - X)$. %If this inequality holds, then we say that $p$ is \emph{locally negamodular} at $X, Y$.
\end{enumerate}
\end{enumerate}
    \item The coverage function $b_{(H, w)}:2^V\rightarrow \Z_{\ge 0}$ is defined by $b_{(H,w)}(X):=\sum_{e\in B_H(X)}w(e)$ for every $X\subseteq V$, where $B_H(X):=\{e\in E: e\cap X\neq \emptyset\}$ for every $X\subseteq V$. 
    \item The hypergraph $(H, w)$ \emph{weakly covers} the function $p$ if $b_{(H, w)}(X)\ge p(X)$ for every $X\subseteq V$. 
    \item The hypergraph $(H, w)$ \emph{strongly covers} the function $p$ if $d_{(H, w)}(X)\ge p(X)$ for every $X\subseteq V$. 
\end{enumerate}
\end{definition}

If a hypergraph $(H, w)$ strongly covers a function $p:2^V\rightarrow \Z$, then it also weakly covers the function $p$. However, the converse is false -- i.e., a weak cover is not necessarily a strong cover\footnote{For example, consider the function $p: 2^V\rightarrow \Z$ defined by $p(X)\coloneqq 1$ for every non-empty proper subset $X\subsetneq V$ and $p(\emptyset)\coloneqq p(V)\coloneqq 0$, and the hypergraph $(H=(V, E\coloneqq \{\{u\}: u\in V\}), w: E\rightarrow \{1\})$.}.  
Bern\'{a}th and Kir\'{a}ly \cite{Bernath-Kiraly} showed 
that a weak cover of a \emph{symmetric skew-supermodular} function can be converted to a strong cover of the same function by repeated \emph{merging of disjoint hyperedges}. 
%a fundamental result that allows one to convert a weak cover to a strong cover via repeated \emph{merging operations}. %(which will be defined shortly).
We recall their definition of the merging operation, discuss their result, and its significance now. 
\begin{restatable}{definition}{defHyperedgeMerge}\label{def:WeakToStrongCover:hyperedge-merge}
%For a hyperedge-weighted hypergraph $(H=(V, E), w: E\rightarrow \Z_+)$, 
Let $(H=(V, E), w: E\rightarrow \Z_+)$ be a hypergraph. 
We use $H_w$ to denote the unweighted multi-hypergraph over vertex set $V$ containing $w(e)$ copies of every hyperedge $e\in E$.  
By \emph{merging two disjoint hyperedges} of $H_w$, we refer to the operation of replacing them by their union in $H_w$. 
We will say that a hypergraph $(G=(V, E_G), c: E_G\rightarrow \Z_+)$ is obtained from $(H, w)$ by \emph{merging hyperedges} if the multi-hypergraph $G_c$ is obtained from the multi-hypergraph $H_w$ by repeatedly merging two disjoint hyperedges in the current hypergraph. 
\end{restatable}
Bern\'{a}th and Kir\'{a}ly showed the following result: 
%Theorem \ref{thm:bernath-kiraly} stated below. Th
\begin{theorem}[Bern\'{a}th and Kir\'{a}ly \cite{Bernath-Kiraly}]\label{thm:bernath-kiraly}
    Let $(H=(V, E), w: E\rightarrow \Z_+)$ be a hypergraph and $p:2^V\rightarrow \Z$ be a symmetric skew-supermodular function such that 
    %$(H,w)$ weakly covers $p$. Then, there exists a hypergraph $(G^*=(V, E^*),w^*: E^*\rightarrow \Z_+)$ obtained by merging hyperedges of the hypergraph $(H,w)$ such that $(G,c)$ strongly covers $p$. 
$b_{(H, w)}(X) \geq p(X)$ for every $X \subseteq V$. 
Then, there exists a hypergraph $\left(H^* = (V, E^*), w^*:E^* \rightarrow\Z_+\right)$ 
such that
\begin{enumerate}[label=$(\arabic*)$]
    \item $d_{(H^*, w^*)}(X) \geq p(X)$ for every $X\subseteq V$ and 
    \item the hypergraph $(H^*,w^*)$ is obtained by merging hyperedges of the hypergraph $(H, w)$.
\end{enumerate}
\end{theorem}

We observe that Theorem \ref{thm:bernath-kiraly} can be used to prove the existential version of Theorem \ref{thm:complete-splitting-off}: namely, for every hypergraph $(G=(V, E), w_G: E\rightarrow \Z_+)$ and a vertex $s\in V$, \emph{there exists} a local connectivity preserving complete h-splitting-off at $s$ from $(G,w_G)$. This can be shown by setting up the hypergraph $(H,w)$ and the function $p$ suitably based on $(G,w_G)$ and using Theorem \ref{thm:bernath-kiraly} (see the first two paragraphs of the proof of Theorem \ref{thm:hypergraphs-splitting-off-stronger-form} in Section \ref{sec:hypergraph-splitting-off-proof}). 
We emphasize that this conclusion regarding hypergraph splitting-off from Bern\'{a}th and Kir\'{a}ly's result was not known before in the literature and is one of our contributions. 
\begin{remark}
    We were also able to prove the existential version of Theorem \ref{thm:complete-splitting-off} using \emph{element-connectivity preserving reduction operations} \cite{BCCKM-private} (see \cite{CK14} for the definition of element-connectivity and the notion of element-connectivity preserving reduction operations) -- we omit the details of this alternate proof in the interests of brevity. The alternate proof does not seem to be helpful for the purposes of a strongly polynomial time algorithm. In fact, it remains open to design a strongly polynomial-time algorithm to perform \emph{complete} element-connectivity preserving reduction operations in the weighted setting \cite{CK14}. 
\end{remark}

We recall that existence of a local connectivity preserving complete h-splitting-off at a vertex from a hypergraph does not immediately imply a polynomial-time algorithm---see the example in Remark \ref{remark:exp-sized-conn-preserving-splitting-off}. 
%algorithm. In particular, there exists a local connectivity preserving complete h-splitting-off at a vertex $s$ from a hypergraph $(G,w_G)$ that contains $2^{O(n)}$ hyperedges although $G$ has only $O(n)$ hyperedges, where $n$ is the number of vertices in $G$ (see the discussion after Theorem \ref{thm:complete-splitting-off} for such an example). 
However, the above-mentioned proof of 
existence of a local-connectivity preserving complete splitting-off at an arbitrary vertex from a hypergraph (i.e., existential version of Theorem \ref{thm:complete-splitting-off}) 
via Theorem \ref{thm:bernath-kiraly} suggests a natural approach towards designing a strongly polynomial time algorithm to find a local-connectivity preserving complete splitting-off at a given vertex from a given hypergraph: it suffices to prove a constructive version of Theorem \ref{thm:bernath-kiraly} via a strongly polynomial-time algorithm. Towards this end, the example in Remark \ref{remark:exp-sized-conn-preserving-splitting-off} suggests a necessary structural step towards a strongly polynomial-time algorithmic version of Theorem \ref{thm:bernath-kiraly}: we need to show Theorem \ref{thm:bernath-kiraly} with the extra conclusion that the number of additional hyperedges in $H^*$ is polynomial in the number of hyperedges and vertices in $H$. 

\iffalse
We observe that Theorem \ref{thm:bernath-kiraly} can be used to prove the existential version of Theorem \ref{thm:complete-splitting-off}: namely, for every hypergraph $(G=(V, E), w_G: E\rightarrow \Z_+)$ and a vertex $s\in V$, \emph{there exists} a local connectivity preserving complete h-splitting-off at $s$ from $(G,w_G)$. This can be shown by setting up the hypergraph $(H,w)$ and the function $p$ suitably based on $(G,w_G)$ and using Theorem \ref{thm:bernath-kiraly} (see the first two paragraphs of the proof of Theorem \ref{thm:hypergraphs-splitting-off-stronger-form} in Section \ref{sec:hypergraph-splitting-off-proof}). 
We emphasize that this conclusion regarding hypergraph splitting-off from Bern\'{a}th and Kir\'{a}ly's result was not known before in the literature and is one of our contributions. 
\knote{The existential version of Theorem \ref{thm:complete-splitting-off} can also be derived from element-connectivity preserving reduction operation. However, neither of these derivations lead to a strongly polynomial-time algorithm.}
However, we recall that existence does not immediately imply a polynomial-time algorithm---see the example in Remark \ref{remark:exp-sized-conn-preserving-splitting-off}. 
%algorithm. In particular, there exists a local connectivity preserving complete h-splitting-off at a vertex $s$ from a hypergraph $(G,w_G)$ that contains $2^{O(n)}$ hyperedges although $G$ has only $O(n)$ hyperedges, where $n$ is the number of vertices in $G$ (see the discussion after Theorem \ref{thm:complete-splitting-off} for such an example). 
This example suggests a necessary structural step towards a strongly polynomial-time algorithmic version of Theorem \ref{thm:bernath-kiraly}: we need to show Theorem \ref{thm:bernath-kiraly} with the extra conclusion that the number of additional hyperedges in $H^*$ is polynomial in the number of hyperedges and vertices in $H$. 
\fi

Bern\'{a}th and Kir\'{a}ly proved Theorem \ref{thm:bernath-kiraly} in the context of a reduction between certain hypergraph connectivity augmentation problems. For that reduction, the existential version of Theorem \ref{thm:bernath-kiraly} is sufficient. However, for the purposes of our application to hypergraph splitting-off, we need an algorithmic version of Theorem \ref{thm:bernath-kiraly}. 
Bern\'{a}th and Kir\'{a}ly's 
proof of Theorem \ref{thm:bernath-kiraly} is in fact algorithmic, but the run-time of the associated algorithm is not necessarily polynomial.  
Their proof implies that the number of additional hyperedges in the hypergraph $(H^*, c^*)$ returned by their algorithm is at most  $\sum_{e\in E}w(e)$ (i.e., $|E^*|-|E|\le \sum_{e\in E}w(e)$) and the run-time of the algorithm is $O(\sum_{e\in E}(|e|+w(e))$. In particular, their run-time is polynomial only if the input weights are given in unary. Moreover, the exponential-sized hypergraph given in the example in Remark \ref{remark:exp-sized-conn-preserving-splitting-off} could indeed arise as a consequence of their algorithm. 
%\knote{Add a reference to the example that appears after Theorem 1.1.}

We address both the structural and the algorithmic issues mentioned above by proving a stronger algorithmic 
%In order to design a strongly polynomial-time algorithm for local connectivity preserving complete h-splitting-off at a vertex $s$ from $(G,w)$, we prove a stronger algorithmic 
version of Theorem \ref{thm:bernath-kiraly}. 
In order to phrase an algorithmic version of Theorem \ref{thm:bernath-kiraly}, we need suitable access to the function $p$. Bern\'{a}th and Kir\'{a}ly \cite{Bernath-Kiraly} suggested access to a certain function maximization oracle associated with the function $p$ that we describe below. 

\begin{definition}
    Let $p: 2^V\rightarrow \Z$ be a set function.    $\functionMaximizationOracleStrongCover{p}\left(\left(G_0, c_0\right), S_0, T_0\right)$ takes as input a hypergraph $(G_0=(V, E_0), c_0: E_0\rightarrow \mathbb{Z}_{+})$ and disjoint sets $S_0, T_0\subseteq V$, and returns a tuple $(Z, p(Z))$, where $Z$ is an optimum solution to the following problem:
        \begin{align*}
        \max & \left\{p(Z)-d_{(G_0, c_0)}(Z): S_0\subseteq Z\subseteq V-T_0\right\}. \tag{\functionMaximizationOracleStrongCover{p}}\label{tag:maximization-oracle-strong-cover}
        \end{align*}
\end{definition}
For the purposes of our application (namely local connectivity preserving complete h-splitting-off at a vertex from a hypergraph), the above-mentioned function maximization oracle can be implemented to run in strongly polynomial time (see Lemma \ref{lemma:helper-for-p-max-oracle}). Using the above mentioned oracle, we prove the following algorithmic version of Theorem \ref{thm:bernath-kiraly}.

\begin{restatable}{theorem}{thmWeakToStrongCover}\label{thm:WeakToStrongCover:main}
Let $(H=(V, E), w: E\rightarrow \Z_+)$ be a hypergraph and $p:2^V\rightarrow \Z$ be a symmetric skew-supermodular function such that 
%$(H,w)$ weakly covers $p$. 
%Then, there exists a hypergraph $(G^*=(V, E^*),w^*: E^*\rightarrow \Z_+)$ obtained by merging hyperedges of the hypergraph $(H,w)$ such that $(G,c)$ strongly covers $p$ and \[|E^*|-|E|=O(|V|).\]
%Let $p:2^V\rightarrow\Z$ be a symmetric skew-supermodular function. Let $\left(H = (V, E), w:E \rightarrow\Z_+\right)$ be a hypergraph such that 
$b_{(H, w)}(X) \geq p(X)$ for every $X \subseteq V$. 
Then, there exists a hypergraph \linebreak $\left(H^* = (V, E^*), w^*:E^* \rightarrow\Z_+\right)$ 
such that
\begin{enumerate}[label=$(\arabic*)$, ref=(\arabic*)]
    \item \label{thm:WeakToStrongCover:main:(1)} $d_{(H^*, w^*)}(X) \geq p(X)$ for every $X\subseteq V$,
    \item \label{thm:WeakToStrongCover:main:(2)} the hypergraph $(H^*,w^*)$ is obtained by merging hyperedges of the hypergraph $(H, w)$, and
    \item \label{thm:WeakToStrongCover:main:(3)} $|E^*| - |E| = O(|V|)$.
\end{enumerate}
%Furthermore, given a hypergraph $(H=(V, E), w: E\rightarrow \Z_+)$ and access to \functionMaximizationOracleStrongCover{p} of a symmetric skew-supermodular function $p: 2^V\rightarrow \Z$ where $b_{(H, w)}(X)\ge p(X)$ for every $X\subseteq V$, there exists an algorithm that runs in time $\poly(|V|, |E|)$ using $\poly(|V|)$ queries to \functionMaximizationOracleStrongCover{p} to return a hypergraph $\left(G^* = (V, E^*), w^*:E^* \rightarrow\Z_+\right)$ satisfying the above three properties. 
Furthermore, given a hypergraph $(H=(V, E), w: E\rightarrow \Z_+)$ and access to \functionMaximizationOracleStrongCover{p} of a symmetric skew-supermodular function $p: 2^V\rightarrow \Z$ where $b_{(H, w)}(X)\ge p(X)$ for every $X\subseteq V$, there exists an algorithm that runs in $O(|V|^3(|V| + |E|)^2)$ time using $O(|V|^3(|V|+|E|))$ queries to \functionMaximizationOracleStrongCover{p} and returns a hypergraph $\left(H^* = (V, E^*), w^*:E^* \rightarrow\Z_+\right)$ satisfying the above three properties. The run-time includes the time to construct the hypergraphs used as input to the queries to \functionMaximizationOracleStrongCover{p}.
Moreover, for each query to \functionMaximizationOracleStrongCover{p}, the hypergraph $(G_0, c_0)$ used as input to the query has $O(|V|)$ vertices and $O(|V|+|E|)$ hyperedges.
\end{restatable}

%Theorem \ref{thm:complete-splitting-off} follows immediately from Theorem \ref{thm:WeakToStrongCover:main} (see Section \ref{sec:hypergraph-splitting-off-proof}). 
%We emphasize that the run-time in Theorem \ref{thm:WeakToStrongCover:main} includes the time to construct the hypergraphs $(G_0, c_0)$ that are used as inputs to \functionMaximizationOracleStrongCover{p}. 
\Cref{thm:WeakToStrongCover:main} is a strengthening of Theorem \ref{thm:bernath-kiraly} in two ways. Firstly, our theorem shows the existence of a hypergraph that not only satisfies properties \ref{thm:WeakToStrongCover:main:(1)} and \ref{thm:WeakToStrongCover:main:(2)}, but also satisfies property \ref{thm:WeakToStrongCover:main:(3)} -- i.e., the number of additional hyperedges in the returned hypergraph is \emph{linear} in the size of the vertex set. Secondly, our \Cref{thm:WeakToStrongCover:main} shows the existence of a strongly polynomial-time algorithm that returns a hypergraph satisfying the three properties. 
%Although Bern\'{a}th and Kir\'{a}ly's result is existential, their proof is algorithmic. However, their proof is based on induction on $\sum_{e\in E}w(e)$ and consequently, the straightforward bound on the number of additional hyperedges as well as the run-time is $O(\sum_{e\in E}w(e))$. Thus, the best run-time bound known for their algorithm was only pseudo-polynomial. 
Our main contribution is modifying Bern\'{a}th and Kir\'{a}ly's algorithm and analyzing the modified algorithm to bound the number of additional hyperedges and the run-time. 
%Our main contribution is analyzing an \emph{adaptation} of Bern\'{a}th and Kir\'{a}ly's algorithm to bound the number of additional hyperedges and the run-time. 
We mention that property \ref{thm:WeakToStrongCover:main:(3)} cannot be tightened to guarantee that $|E^*-E|=O(|V|)$ -- we were able to construct an example where 
%the only strong cover $(H^*=(V, E^*), w^*: E^*\rightarrow \Z_+)$ of a function $p: 2^V\rightarrow \Z$ that is obtained by merging hyperedges of a weak cover hypergraph $(H=(V, E), w: E\rightarrow \Z_+)$ of the function $p$ is such that 
$|E^*-E|=\Omega(|V|^2)$ (see \Cref{sec:appendix:tight-example-weak-to-strong}). 

%Our \Cref{thm:WeakToStrongCover:main} shows that there exists a hypergraph $(G, c)$ obtained from $(H, w)$ by merging hyperedges such that $(G, c)$ strongly covers $p$ and the number of additional hyperedges in $G$ compared to $H$ is $O(|V|)$ (i.e., $|E_G|-|E_H|=O(|V|)$); moreover, our proof is constructive and leads to a strongly polynomial-time algorithm to find such a strong cover $(G, c)$ of $p$ assuming access to the function $p$ via \functionMaximizationOracleStrongCover{p}. This result leads to a strongly polynomial-time algorithm for hypergraph splitting-off (see \Cref{sec:splitting-off}). 

\Cref{thm:WeakToStrongCover:main} immediately leads to  a proof of \Cref{thm:complete-splitting-off} (see Theorem \ref{thm:hypergraphs-splitting-off-stronger-form} and its proof in Section \ref{sec:hypergraph-splitting-off-proof}). Instead of using \Cref{thm:WeakToStrongCover:main} as a black-box, if we delve into the proof of it in the context of the proof of \Cref{thm:complete-splitting-off}, we obtain the following theorem: 

%In the proof of \Cref{thm:complete-splitting-off}, we used \Cref{thm:WeakToStrongCover:main} as a black-box. Instead, delving into the proof of \Cref{thm:WeakToStrongCover:main} in the context of the proof of \Cref{thm:complete-splitting-off} reveals that the following theorem also holds. 
\begin{theorem}\label{thm:splitting-off-or-trim-to-preserve-connectivity}
    Let $(G=(V, E), w: E\rightarrow \Z_+)$ be a hypergraph and $s\in V$. Then, there exists a hypergraph $(G'=(V, E'), w': E'\rightarrow \Z_+)$ obtained 
    by applying a h-splitting-off operation at $s$ from $(G, w)$ 
    %from $(G, w)$ by a h-splitting-off operation 
    such that $\lambda_{(G', w')}(u, v)=\lambda_{(G, w)}(u, v)$ for every distinct $u, v\in V\setminus \{s\}$. 
\end{theorem}
We omit the proof of \Cref{thm:splitting-off-or-trim-to-preserve-connectivity} in the interests of brevity. 
\Cref{thm:splitting-off-or-trim-to-preserve-connectivity} closely resembles the \emph{existential edge splitting-off} results of Lov\'{a}sz \cite{Lov74, Lovasz-problems-book} and Mader \cite{Mad78} for graphs. Lov\'{a}sz's and Mader's  existential edge splitting-off results for graphs are important since they have been used to simplify the proofs of fundamental results in graph theory---e.g., Nash-Williams' Strong Orientation Theorem. 
On the other hand, \Cref{thm:splitting-off-or-trim-to-preserve-connectivity} does not immediately imply a strongly polynomial-time algorithm for finding a local connectivity preserving complete h-splitting off at a vertex from a given weighted hypergraph. So, \Cref{thm:complete-splitting-off} may be useful in algorithmic contexts while \Cref{thm:splitting-off-or-trim-to-preserve-connectivity} may be useful in graph-theoretical contexts. 

\subsection{Proof Technique for \Cref{thm:WeakToStrongCover:main}}
In this section, we describe our proof technique for the existential result in \Cref{thm:WeakToStrongCover:main}. 
The algorithmic results in that theorem follow from the existential result using standard algorithmic tools for submodular functions, so we focus only on outlining a proof of the existential result. 
Let $(H=(V, E), w: E\rightarrow \Z_+)$ be a hypergraph and $p: 2^V\rightarrow \Z_+$ be a symmetric skew-supermodular function such that $(H, w)$ weakly covers the function $p$. Our goal is to show that there exists a hypergraph $(H^*=(V, E^*), w^*: E\rightarrow \Z_+)$ such that 
\begin{enumerate}[label=(\arabic*)]
    \item \hypertarget{WeakToStrong(1)}{$(H^*, w^*)$ strongly covers the function $p$},
    \item \hypertarget{WeakToStrong(2)}{$(H^*, w^*)$ is obtained by merging hyperedges of the hypergraph $(H, w)$}, and
    \item \hypertarget{WeakToStrong(3)}{$|E^*|-|E|=O(|V|)$}. 
\end{enumerate}
% \hypertarget{WeakToStrong(1)}{(1) $(H^*, w^*)$ strongly covers the function $p$}, \hypertarget{WeakToStrong(2)}{(2) $(H^*, w^*)$ is obtained by merging hyperedges of the hypergraph $(H, w)$}, and \hypertarget{WeakToStrong(3)}{(3) $|E^*|-|E|=O(|V|)$}. 

\paragraph{Preliminaries.} We define a set $X\subseteq V$ to be $(p, H, w)$-tight if $b_{(H, w)}(X) = p(X)$. For a function $p$ and hypergraph $(H, w)$, let $T_{p, H, w}$ denote the family of $(p,H, w)$-tight sets and let $\mathcal{T}_{p, H, w}$ be the family of inclusionwise maximal sets in $T_{p, H, w}$. We will need the following two operations: 
\begin{enumerate}[label=(\roman*)]
    \item For hyperedges $e, f\in E$ and a positive integer  $\alpha \le \min\{w(e), w(f)\}$, the operation \merge$((H, w), e, f, \alpha))$ returns the hypergraph obtained from $(H, w)$ by decreasing the weight of hyperedges $e$ and $f$ by $\alpha$ and increasing the weight of the hyperedge $e\cup f$ by $\alpha$. All hyperedges with zero weight are discarded.
    \item For a hyperedge $e\in E$ and a positive integer $\alpha\le w(e)$, the operation \reduce$((H, w), e, \alpha)$ returns the hypergraph obtained by decreasing the weight of the hyperedge $e$ by $\alpha$. All hyperedges with zero weight are discarded.
\end{enumerate}
 
%The sets in $\mathcal{T}_{p, H, w}$ can be shown to be pairwise disjoint using the hypotheses that $p$ is skew-supermodular and $(H,w)$ weakly covers $p$ \shubhang{The family $\maximalTightSetFamily{p, H, w}$ being disjoint requires a stronger hypothesis (see \Cref{lem:weak-to-strong:uncrossing:maximal-tight-set-family-disjoint}). This is the main motivation for enforcing that every hyperedge is contained in a tight set and every vertex is contained in a hyperedge during each recursive call of the algorithm. Consequently, this is also why we immediately do reduction and contraction steps after the merge step. This insight could also count as one of our key technical contributions.}. 

%For two hypergraphs $(G=(V, E_G), c_G:E_G\rightarrow \Z_+)$ and $(H=(V, E_H), w_H: E_H\rightarrow \Z_+)$ on the same vertex set $V$, we define the hypergraph $(G+H=(V, E_{G+H}), c_G+w_H)$ as the hypergraph with vertex set $V$ and hyperedge set $E_{G+H}\coloneqq E_G\cup E_H$ with the weight of every hyperedge $e\in E_G\cap E_H$ being $c_G(e) + w_H(e)$, the weight of every hyperedge $e\in E_G\setminus E_H$ being $c_G(e)$,  and the weight of every hyperedge $e\in E_H\setminus E_G$ being $w_H(e)$. 


\paragraph{Algorithm of \cite{Bernath-Kiraly}.} Our proof of the existential result builds on the techniques of Bern\'{a}th and Kir\'{a}ly \cite{Bernath-Kiraly} who proved the existence of a hypergraph $(H^*=(V, E^*), w^*: E^*\rightarrow \Z_+)$ satisfying properties \hyperlink{WeakToStrong(1)}{(1)} and \hyperlink{WeakToStrong(2)}{(2)}, so we briefly recall their techniques. 
We present the algorithmic version of their proof since it will be useful for our purposes. 

The proof in \cite{Bernath-Kiraly} is inductive, and consequently, the algorithm implicit in their proof is recursive. The algorithm takes as input a hypergraph $((H=(V, E), w: E\rightarrow \Z_+)$ and a symmetric skew-supermodular function $p: 2^V\rightarrow \Z$ such that the hypergraph $(H,w)$ weakly covers the function $p$. 
%The algorithm returns a hypergraph $(H^*, w^*)$ that satisfies properties (1) and (2).
If $w(E) = 0$, then the algorithm is in its base case and returns the empty hypergraph. 
Otherwise, $w(E) > 0$; the algorithm chooses an arbitrary hyperedge $e \in E$ and defines hypergraphs $(H_0, w_0)$ and $(H', w')$ and the function $p'$ by considering two cases.
First, suppose that the hyperedge $e$ is not contained in any set of the family $\maximalTightSetFamily{p, H, w}$. 
In this case, the algorithm defines $(H_0, w_0)$ to be the hypergraph on vertex set $V$ consisting of a single hyperedge $e$ with $w_0(e) = 1$, constructs the hypergraph $(H', w'):=\reduce((H, w), e, 1)$, and defines the function $p' := p - d_{(H_0, w_0)}$. Second, suppose that the hyperedge $e$ is contained in some set $X \in \mathcal{T}_{p, H, w}$. It can be shown that there exists a hyperedge $f \in E$ such that $f \subseteq V - X$. In this case, the algorithm defines $(H_0, w_0)$ to be the empty hypergraph on vertex set $V$, constructs the hypergraph $(H', w'):=\merge((H, w), e, f, 1)$, and defines the function $p' := p$. In both cases, the algorithm recurses on the inputs $(H',w')$ and $p'$ to obtain a hypergraph $(H^*_0, w^*_0)$ and returns the hypergraph $(H^*, w^*) = (H^*_0 + H_0, w^*_0 + w_0)$. 
Here, the hyperedges of $H^*$ are the union of the hyperedges of $H^*_0$ and $H_0$ with the weight $w^*(e)$ of a hyperedge $e$ being the sum of the weights $w^*_0(e)+w_0(e)$ if the hyperedge $e$ is present in both $H^*_0$ and $H_0$, being $w^*_0(e)$ if the hyperedge $e$ is present only in $H^*_0$, and being $w_0(e)$ if the hyperedge $e$ is present only in $H_0$. 
\iffalse
\knote{Edit this to be consistent with BK's proof.}\shubhang{Rest of this paragraph can be removed.}
First, suppose that there exists a hyperedge $e\in E$ that is not contained in any set of the family $\maximalTightSetFamily{p, H, w}$. 
In this case, the algorithm defines $(H_0, w_0)$ to be the hypergraph on vertex set $V$ consisting of a single hyperedge $e$ with $w_0(e) = 1$,
constructs the hypergraph $(H', w'):=\reduce((H, w), e, 1)$, and defines the function $p' := p - d_{(H_0, w_0)}$.
Second, suppose that every hyperedge $e\in E$ is contained in some set of the family $\mathcal{T}_{p, H, w}$. It can be shown that there exists a pair of disjoint hyperedges $e, f\in E$ contained in distinct sets of the family $\maximalTightSetFamily{p, H, w}$.
In this case, the algorithm defines $(H_0, w_0)$ to be the empty hypergraph on vertex set $V$, constructs the hypergraph $(H', w'):=\merge((H, w), e, f, 1)$, and defines the function $p' := p$.
In both cases, the algorithm recurses on the inputs $(H',w')$ and $p'$ to obtain a hypergraph $(H^*_0, w^*_0)$ and returns the hypergraph $(H^*, w^*) = (H^*_0 + H_0, w^*_0 + w_0)$. 
Here, the hyperedges of $H^*$ are the union of the hyperedges of $H^*_0$ and $H_0$ with the weight $w^*$ of a hyperedge $e$ being the sum of the weights $w^*_0(e)+w_0(e)$ if the hyperedge $e$ is present in both $H^*_0$ and $H_0$, being $w^*_0(e)$ if the hyperedge $e$ is present only in $H^*_0$, and being $w_0(e)$ if the hyperedge $e$ is present only in $H_0$. 
\fi

 We note that $w'(E') = w(E) - 1$. Furthermore, it can be shown that the function $p'$ is symmetric skew-supermodular and the hypergraph $(H', w')$ weakly covers the function $p'$. Consequently, by induction on $w(E)$, the algorithm can be shown to terminate in $w(E)$ recursive calls and returns a hypergraph satisfying properties \hyperlink{WeakToStrong(1)}{(1)} and \hyperlink{WeakToStrong(2)}{(2)}. Moreover, the number of additional hyperedges in the returned hypergraph is at most the
number of recursive calls where the \merge operation is performed, which is also at most $w(E)$. 
%\knote{Add a reference to the example that appears after Theorem 1.1.}
Thus, in order
to reduce the number of additional hyperedges and to design a strongly polynomial-time algorithm, our goal is to
reduce the recursion depth of the algorithm. We emphasize that the recursion depth of Bern\'{a}th and Kir\'{a}ly's algorithm could indeed be exponential (the exponential sized example mentioned in Remark \ref{remark:exp-sized-conn-preserving-splitting-off} could arise in the execution of their algorithm), so we need to necessarily modify their algorithm. 

\paragraph{Preprocessing for Additional Structure.}
%Our first technical insight is that by performing a preprocessing step (see Section \ref{sec:WeakToStrongCover}), we can assume that the inputs to our algorithm satisfy additional conditions. \shubhang{Furthermore, appropriate modifications to Bern\'{a}th and Kir\'{a}ly's algorithm ensures that these conditions hold during every recursive call of our recursive algorithm described in the paragraph below.}
%We now describe our adaptation of the above-mentioned Bern\'{a}th and Kir\'{a}ly's algorithm to reduce the recursion depth. 
Similar to Bern\'{a}th and Kir\'{a}ly's algorithm, our algorithm also takes as input a hypergraph $(H=(V, E), w: E\rightarrow \Z_+)$ and a symmetric skew-supermodular function $p: 2^V\rightarrow \Z$ such that the hypergraph $(H,w)$ weakly covers the function $p$. 
However, unlike Bern\'{a}th and Kir\'{a}ly's algorithm, our algorithm performs a preprocessing step so that the inputs $(H,w)$ and $p$ satisfy the following two additional conditions: 
\begin{enumerate}[label=(\alph*)]
    \item \hypertarget{WeakToStrong(a)}{every hyperedge $e\in E$ is contained in some set of $\mathcal{T}_{p, H, w}$} and 
    \item \hypertarget{WeakToStrong(b)}{the degree of every vertex in $(H, w)$ is non-zero.}
\end{enumerate}
%These two additional conditions will be inductively maintained by our algorithm. 
As a consequence of these additional conditions, the family $\maximalTightSetFamily{p, H, w}$ will be a \emph{disjoint family},  
%during every recursive call of our algorithm, 
a property that we leverage heavily throughout our analysis.
% \hypertarget{WeakToStrong(a)}{(a) every hyperedge $e\in E$ is contained in some set of $\mathcal{T}_{p, H, w}$} and \hypertarget{WeakToStrong(b)}{(b) the degree of every vertex in $(H, w)$ is non-zero} -- we show that these conditions are without loss of generality since an arbitrary instance can be reduced to an instance satisfying these conditions (see Section \ref{sec:WeakToStrongCover}).
Furthermore, we modify Bern\'{a}th and Kir\'{a}ly's algorithm to ensure that these conditions hold during every recursive call.

\paragraph{Our Algorithm.}
We now describe our modification of the above-mentioned Bern\'{a}th and Kir\'{a}ly's algorithm to reduce the recursion depth. Our algorithm is also recursive and its 
%, and returns a hypergraph $(H^*, w^*)$ that satisfies properties (1), (2) and (3).
base case is the same as that of Bern\'{a}th and Kir\'{a}ly's algorithm (i.e., $w(E)=0$). During recursive cases (i.e., if $w(E) > 0$), instead of performing one of the two (i.e., \merge or  \reduce) operations, our algorithm performs both operations in a sequential fashion. In particular, we find a pair of disjoint hyperedges $e, f\in E$ contained in distinct sets of $\maximalTightSetFamily{p, H, w}$ (such a pair exists by condition \hyperlink{WeakToStrong(a)}{(a)} and the arguments of Bern\'{a}th and Kir\'{a}ly mentioned above). Next, 
instead of performing one \merge operation (as was done by Bern\'{a}th and Kir\'{a}ly's algorithm), we perform as many  \merge operations as possible %\shubhang{(denoted by $\alpha^M \in \Z_+$)} 
using the hyperedges $e$ and $f$. 
Formally, let $$\alpha^M := \max\left\{\alpha \in \Z_+ : \text{hypergraph returned by } \merge((H, w), e, f, \alpha) \text{ weakly covers } p\right\}.$$
We let $(H^M, w^M):=\textsc{Merge}((H, w), e, f, \alpha^M)$ and $p^M:=p$.
%We describe this formally now: we obtain $(H^M, w^M):=\textsc{Merge}((H, w), e, f, \alpha^M)$ and $p^M:=p$, where $\alpha^M\le \min\{w(e), w(f)\}$ is the largest integer such that $(H^M, w^M)$ weakly covers the function $p^M$. 
Next, instead of recursing on $((H^M, w^M), p^M)$ (as was done by Bern\'{a}th and Kir\'{a}ly's algorithm), we perform as many \reduce operations as possible 
%\shubhang{(denoted by $\alpha^R \in \Z_{\geq 0}$)} 
on the newly created hyperedge $e\cup f$.
Formally, let $$\alpha^R := \max\left\{\alpha \in \Z_{\ge 0}: \text{hypergraph returned by } \reduce((H^M, w^M), e \cup f, \alpha) \text{ weakly covers } (p^M - d_{(H_0^\alpha, w_0^\alpha)})\right\},$$
%where $(H_0^{\alpha}:=(V, E_0^{\alpha}:=\{e\cup f\}),w_0^{\alpha}: E_0^{\alpha}\rightarrow \{\alpha\})$ 
where $(H_0^{\alpha}, w_0^{\alpha})$ denotes
the hypergraph on vertex set $V$ consisting of a single hyperedge $e\cup f$ with $w_0^{\alpha}(e\cup f) = \alpha$. 
We construct the hypergraph 
$(H_0, w_0):=(H_0^{\alpha^R}, w_0^{\alpha^R})$, the hypergraph $(H^R, w^R):=\reduce((H^M, w^M), e\cup f, \alpha^R)$ and define the function $p^R:=p^M - d_{(H_0^{\alpha^R}, w_0^{\alpha^R})}$.
%We describe this formally now: 
%for $\alpha\le w^M(e\cup f)$, let $(H_0^{\alpha}:=(V, E_0^{\alpha}:=\{e\cup f\}),w_0^{\alpha}: E_0^{\alpha}\rightarrow \{\alpha\})$. 
%Let $(H^R, w^R):=\textsc{Reduce}((H^M, w^M), e\cup f, \alpha^R)$ and $p^R:=p^M-d_{(H_0^{\alpha^R},w_0^{\alpha^R})}$, where $\alpha^R\le w^M(e\cup f)$ is the largest integer such that $(H^R, w^R)$ weakly covers the function $p^R$. 
%and $(H^R, w^R):=(H^R_{\alpha_R}, w^R_{\alpha_R})$, 
%and define the function $p^R:=p^R_{\alpha_R}$. 
This immediate reduce step ensures that the hypergraph $(H^R, w^R)$ and the function $p^R$ satisfy condition \hyperlink{WeakToStrong(a)}{(a)} -- i.e., every hyperedge in $(H^R, w^R)$ is contained in some set of $\mathcal{T}_{p^R, H^R, w^R}$. 
Finally, we compute sets $\zeros:=\{u\in V: b_{(H^R, w^R)}(u)=0\}$ and $V':=V-\zeros$, hypergraph $(H':=(V', E':=E^R), w':=w^R)$, and define the function $p':2^{V'}\rightarrow \Z$ by $p'(X):=\max\{p(X\cup R): R\subseteq \zeros\}$ for every $X\subseteq V'$ -- this final step can be viewed as a clean up step since it gets rid of vertices that are not incident to any hyperedges (and revises the function $p$ appropriately). 
This clean up step ensures that the hypergraph $(H', w')$ satisfies condition \hyperlink{WeakToStrong(b)}{(b)} -- i.e., the degree of every vertex in $(H', w')$ is non-zero. 
It can be shown that the function $p'$ is symmetric skew-supermodular and the hypergraph $(H', w')$ weakly covers the function $p'$. Furthermore, the function $p'$ and hypergraph $(H', w')$ satisfy conditions \hyperlink{WeakToStrong(a)}{(a)} and \hyperlink{WeakToStrong(b)}{(b)}. 
We recursively call the algorithm on input $((H', w'), p')$ to obtain a hypergraph $(H_0^*, w_0^*)$. We obtain the hypergraph $(G, c)$ from $(H_0^*, w_0^*)$ by adding the vertices $\zeros$ and return the hypergraph $(G+H_0, c+w_0)$.  
By induction on the total weight of hyperedges in the input hypergraph, it can be shown that our algorithm returns a hypergraph satisfying properties \hyperlink{WeakToStrong(1)}{(1)} and \hyperlink{WeakToStrong(2)}{(2)} and also terminates within a finite number of recursive calls. 

\paragraph{Recursion Depth and Potential Functions.} We now sketch our proof to show that the recursion depth of our algorithm is $|E|+O(|V|)$. We note that this also bounds the number of additional hyperedges in the hypergraph returned by the algorithm, and consequently this hypergraph also satisfies property \hyperlink{WeakToStrong(3)}{(3)}. Let $\ell$ be the number of recursive calls made by the algorithm on the input instance $((H, w), p)$. 
We partition 
%\shubhang{The proof doesn't exactly use a partition (i.e. there can be an overlap in the two sets of recursive calls that are analyzed separately in the proof). This should be fine though.} 
the set $[\ell]$ of recursive calls into two parts: let $P_1\subseteq [\ell]$ be the set of recursive calls during which 
the merged hyperedge $e\cup f$ survives in the hypergraph $(H', w')$ that is input to the subsequent recursive call and 
$P_2\subseteq [\ell]$ be the recursive calls during which 
the merged hyperedge $e\cup f$ does not survive in the hypergraph $(H', w')$ that is input to the subsequent recursive call. 
We note that $[\ell]=P_1\uplus P_2$. We bound $|P_1|$ and $|P_2|$ separately using certain carefully designed potential functions.

First, we show that $|P_1|=O(|V|)$ as follows: for $i\in [\ell]$, consider the maximal tight set family $\mathcal{T}_i=\mathcal{T}_{p_i, H_i, w_i}$ where $((H_i, w_i), p_i)$ is the input to the $i^{th}$ recursive call. Also, let $\mathcal{T}_{\le 1}:=\mathcal{T}_1$ and $\mathcal{T}_{\le i}:=\mathcal{T}_{i}\cup \{X\cap V_i: X\in \mathcal{T}_{\le i-1}\}$ for integers $i$ where $2 \le i\le \ell$ and $V_i$ is the ground set of the input to the $i^{th}$ recursive call. Thus, $\calT_{\leq i}$ is the \emph{projection} of all the maximal tight sets encountered in the first $i$ recursive calls of the algorithm onto the ground set of the inputs at the $i^{th}$ recursive call. We show that $\mathcal{T}_{\le i}$ is laminar for every $i\in [\ell]$ (\Cref{lem:WeakToStrong:properties-of-cumulative-tight-set-families}). 
However, $|\mathcal{T}_{\le i}|$ is not necessarily non-decreasing with $i$ since projection of a set family to a subset could result in the loss of sets from the family. Consequently, $|\mathcal{T}_{\le i}|$ is not suitable as a potential function to measure progress. Instead, we use the potential function $\phi(i):=|\mathcal{T}_{\le i}|+3|\zeros_{\le i-1}|$, where $\zeros_{\le i}$ is the union of the sets $\zeros$ computed up to the $i^{th}$ recursive call. We show that $\phi(i)$ is non-decreasing and strictly increases if $i\in P_1$ (\Cref{claim:WeakToStrong:hyperedge-support-size:alpha^R=weight-of-added-hyperedge} in \Cref{lem:WeakToStrongCover:NumberOfHyperedges}). Consequently, $|P_1|=O(|V|)$.

Secondly, we bound $|P_2|$ as follows. We use a lookahead-potential function: let $\Phi_1(i)$ be the number of recursive calls between $i$ and $\ell$
during which the merged hyperedge $e\cup f$ survives in the hypergraph $(H', w')$ that is input to the subsequent recursive call
and let $\Phi(i):=|E_i|+\Phi_1(i)$, where $E_i$ is the set of hyperedges in the hypergraph $(H_i, w_i)$ input to the $i^{th}$ recursive call. We show that $\Phi(i)$ is non-increasing and strictly decreases if 
$i\in P_2$ 
(\Cref{claim:WeakToStrong:hyperedge-support-size:alpha^R<weight-of-added-hyperedge} in \Cref{lem:WeakToStrongCover:NumberOfHyperedges}). Hence, $|P_2|\le \Phi(1)-\Phi(\ell)\le \Phi(1) \le |E_1|+|P_1|=|E|+O(|V))$, where the last equality is because of the bound on $|P_1|$ from the previous paragraph.  

\begin{remark}
    Our key technical contributions are twofold. Our first key technical contribution is identifying conditions \hyperlink{WeakToStrong(a)}{(a)} and \hyperlink{WeakToStrong(b)}{(b)} under which $\maximalTightSetFamily{p, H, w}$ becomes a disjoint family. We ensure that conditions \hyperlink{WeakToStrong(a)}{(a)} and \hyperlink{WeakToStrong(b)}{(b)} hold in every recursive call by performing \emph{immediate} reduction and clean-up steps in the algorithm. 
    Our second key technical contribution is identifying appropriate potential functions to measure progress of the algorithm. The disjointness of $\maximalTightSetFamily{p, H, w}$ was crucial for identifying the laminar structure of the family of projected maximal tight sets across recursive calls, which was subsequently helpful in bounding the number of recursive calls corresponding to $P_1$. 
    %Our third key technical contribution is bounding the number of recursive calls corresponding to $P_2$ required a lookahead-potential function that relates $|P_2|$ to $|P_1|$. As discussed above, the additive $|E|$ in the recursion depth comes from the bound on $|P_2|$.
    Moreover, we bound the number of recursive calls corresponding to $P_2$ using a lookahead-potential function that relates $|P_2|$ to $|P_1|$. As discussed above, the additive $|E|$ in the recursion depth comes from the bound on $|P_2|$.
\end{remark}